\subsection{Integral coefficient-tail extension groups}

We use the abelian category $\mathrm{gr}S$ of integer-weight-graded
$S$-modules and weight-preserving homomorphisms. The symbol
$\operatorname{Ext}^i_{\mathrm{gr}S}$ denotes Ext in that category.
Put
\[
 S=\Gamma_{\mathbb Z}(T),\quad
 I_N=\bigoplus_{r\geq N}\mathbb ZT^{[r]},\quad B_N=S/I_N
 \qquad(N\geq1).
\]
Our shift convention is $(M\langle j\rangle)^r=M^{r-j}$;
thus $S\langle j\rangle$ has its free generator in weight $j$.
When both source and target have nonnegative weights, the
displayed Ext groups also compute Ext in the category of
nonnegatively weight-graded modules. Indeed a module in that
category has a free graded resolution generated in nonnegative
weights: choose homogeneous generators, then repeat on the
nonnegative-weight kernels. This same resolution is projective
in the full graded category, and the two Hom complexes are
identical. Negative shifts used in the proof are interpreted
in the full integer-graded category before translating back.
Define
\[
 \widehat{\mathbb Z}=\prod_p\mathbb Z_p,
 \qquad \mathbb A_f=\widehat{\mathbb Z}\otimes_{\mathbb Z}\mathbb Q.
\]
These are algebraic modules and rings in the following arguments;
no topology or completed tensor product is inserted.

\paragraph{AT01 (Proved; coefficient arithmetic sequence).}
There is a pure exact sequence of abelian groups
\[
 0\longrightarrow\mathbb Z\longrightarrow
 \mathbb Q\oplus\widehat{\mathbb Z}
 \xrightarrow{(a,b)\mapsto b-a}\mathbb A_f\longrightarrow0.
\]
Its tensor product with $S$ is an exact sequence of graded
$S$-modules. In the category of graded modules,
$S\otimes\widehat{\mathbb Z}=\prod_p(S\otimes\mathbb Z_p)$.

\begin{proof}
The embedding of $\widehat{\mathbb Z}$ into $\prod_p\mathbb Q_p$
extends injectively after tensoring with $\mathbb Q$, since an
element has one common integer denominator. A rational number
integral at every prime is an integer, proving the kernel
assertion. Write an element of $\mathbb A_f$ as $b/D$, with
$D>0$ and $b\in\widehat{\mathbb Z}$. For each prime dividing
$D$, choose an integer $c$ satisfying
$c\equiv b_p\pmod{p^{v_p(D)}\mathbb Z_p}$, using the Chinese
remainder theorem for this finite set of congruences. Then
$b/D-c/D\in\widehat{\mathbb Z}$. The pair
$(-c/D,b/D-c/D)$ maps to $b/D$. This proves surjectivity.
The quotient $\mathbb A_f$ is a rational vector space, hence
flat over $\mathbb Z$, so the sequence is pure. Tensoring
with $S$ gives the claimed exact sequence. The product
identity is degreewise: each $S^r$ is free of rank one over
$\mathbb Z$. It is an identity of graded products; we make
no assertion identifying it with the ungraded product.
\end{proof}

\paragraph{AT02 (Proved; complete calculation of the tail Ext groups).}
For every $N\geq1$ and every integer $j$,
\[
 \operatorname{Hom}_{\mathrm{gr}S}
       (I_N,S\langle j\rangle)\cong S^{-j},
\]
where $S^{-j}=0$ for $j>0$. The isomorphism is restriction
from $\operatorname{Hom}_{\mathrm{gr}S}(S,S\langle j\rangle)$.
Moreover
\[
 \operatorname{Ext}^1_{\mathrm{gr}S}(I_N,S\langle j\rangle)
 \cong
 \begin{cases}
  \mathbb A_f/\mathbb Q,&1\leq j\leq N,\\
  0,&j\leq0\text{ or }j>N,
 \end{cases}
\]
and $\operatorname{Ext}^i_{\mathrm{gr}S}(I_N,S\langle j\rangle)=0$
for every $i\geq2$. Under the displayed identification, restriction
from $I_N$ to $I_M$, $M\geq N\geq j\geq1$, is the identity of
$\mathbb A_f/\mathbb Q$.

\begin{proof}
We first compute Ext from $B_N$. For every prime $p$, the
flat extension $S\to S_p=S\otimes\mathbb Z_p$ gives derived
adjunction
\[
 R\operatorname{Hom}_{\mathrm{gr}S}
       (B_N,S_p\langle j\rangle)
 =R\operatorname{Hom}_{\mathrm{gr}S_p}
       ((B_N)_p,S_p\langle j\rangle).
\]
This equality can be verified with a graded projective
resolution of $B_N$: tensoring it with the flat algebra
$S_p$ remains a projective resolution, and the ordinary Hom
adjunction holds termwise. The source $(B_N)_p$ is bounded
in weight, hence is torsion in the precise sense of
Nagpal--Snowden: every element is killed by all sufficiently
large divided-power basis elements. The target is free of
rank one, and hence finitely presented, over $S_p$.
The base ring $\mathbb Z_p$ is noetherian and complete for
$(p)$, and its classical divided-power sequence has
$\pi_{p^a}=p$ for every $a\geq1$.
Nagpal--Snowden, Theorem~8.3, p.~45, with proof on
pp.~46--48 of arXiv:1606.03431v2, therefore gives zero
for this derived Hom in every cohomological degree
\cite{IC-nagpal-snowden}. This invocation requires no finite
presentation assumption on the source.

Derived Hom into a graded product is the product of the
derived Homs here. Indeed products are exact in graded
module categories, and products of injectives are injective:
applying Hom into a product gives the product of the exact
Hom functors. Using the graded product in AT01, we conclude
\[
 R\operatorname{Hom}_{\mathrm{gr}S}
 (B_N,S\otimes\widehat{\mathbb Z}\langle j\rangle)=0.
\]

Over $\mathbb Q$ write $t=T$. The map
$T^{[r]}\mapsto t^r/r!$ identifies $S_{\mathbb Q}$ with
$\mathbb Q[t]$ as a weight-graded algebra and identifies
$(B_N)_{\mathbb Q}$ with $\mathbb Q[t]/(t^N)$.
Its graded free resolution is
\[
 0\longrightarrow\mathbb Q[t]\langle N\rangle
   \xrightarrow{\ t^N\ }\mathbb Q[t]
   \longrightarrow\mathbb Q[t]/(t^N)\longrightarrow0.
\]
For a rational vector space $V$, viewed as coefficients in
$V[t]\langle j\rangle$, the degree-preserving Hom complex
of this resolution has zero $H^0$; its $H^1$ is $V$
exactly for $1\leq j\leq N$, and is zero otherwise.
To see this without a grading convention implicit in the
calculation, its terms are $V[t]^{-j}$ and $V[t]^{N-j}$.
Multiplication by $t^N$ is an isomorphism when both terms
are present, and precisely the second term is present in
the indicated range. There are no higher terms.

Apply $R\operatorname{Hom}_{\mathrm{gr}S}(B_N,-)$ to the
graded arithmetic sequence of AT01, shifted by $j$.
The completed-product term is zero by the preceding
argument. The rational and adelic terms have just been
computed, with $V=\mathbb Q$ and $V=\mathbb A_f$.
Their map on degree-one Ext is the diagonal inclusion
$\mathbb Q\to\mathbb A_f$. The long exact sequence gives
\[
 \operatorname{Ext}^i_{\mathrm{gr}S}(B_N,S\langle j\rangle)
 =
 \begin{cases}
  \mathbb A_f/\mathbb Q,&i=2,\ 1\leq j\leq N,\\
  0,&\text{otherwise}.
 \end{cases}
\]
Finally apply graded Hom to
$0\to I_N\to S\to B_N\to0$. Since $S$ is free over
itself, its positive Ext groups vanish. The preceding
formula proves the Hom assertion and every stated Ext
assertion for $I_N$.

For the transition map, use a representative $c\in\mathbb A_f$
and the rational multiplier $c\,t^{-j}$ on $I_N$.
This multiplier is defined in all source weights because
$N\geq j$. Restriction to $I_M$ uses the same multiplier.
The arithmetic connecting homomorphism identifies its
extension class with $c$ modulo $\mathbb Q$, proving the
claimed identity on the quotient. The next statement
makes this representative explicit.
\end{proof}

\paragraph{AT03 (Proved; explicit persistent nonsplit tail extensions).}
Fix $1\leq j\leq N$ and $c\in\mathbb A_f$. Define the
weight-preserving $S$-linear map
\[
 \phi_c:I_N\longrightarrow S_{\mathbb A_f}\langle j\rangle,
 \qquad
 \phi_c(T^{[r]})=
 c\frac{(r-j)!}{r!}\,T^{[r-j]}e_j\quad(r\geq N),
\]
where $e_j$ is the shifted free generator. Form the graded
module
\[
 E_{c,N,j}=\{(a,b,z)\in
 S_{\mathbb Q}\langle j\rangle\oplus
 S_{\widehat{\mathbb Z}}\langle j\rangle\oplus I_N:
 b-a=\phi_c(z)\}.
\]
It fits into the explicit exact sequence
\[
 0\longrightarrow S\langle j\rangle
 \xrightarrow{s\mapsto(s,s,0)}E_{c,N,j}
 \xrightarrow{(a,b,z)\mapsto z}I_N\longrightarrow0.
\]
The sequence splits by a graded $S$-linear map exactly when
$c\in\mathbb Q$. Its pullback along $I_M\hookrightarrow I_N$
is $E_{c,M,j}$, so nonsplitting persists for every $M\geq N$.
For example, the coordinate idempotent $c=e_p$ of
$\widehat{\mathbb Z}$ gives a nonsplit extension.

\begin{proof}
The formula for $\phi_c$ is multiplication by $c t^{-j}$
under the explicitly given rational polynomial comparison;
it is therefore $S$-linear and preserves weight.
Exactness is the pullback of AT01 in each weight.
A section is equivalent to a pair of graded $S$-linear maps
from $I_N$ to the rational and completed-product terms whose
difference is $\phi_c$. The completed-product Hom group is
zero: at every prime the exact sequence
$0\to(I_N)_p\to S_p\to(B_N)_p\to0$, together with the
vanishing furnished by Theorem~8.3, identifies this Hom
group with $(S_p\langle j\rangle)^0=0$.
The rational Hom group consists exactly of the multipliers
$a t^{-j}$, $a\in\mathbb Q$, since the rationalized tail
is the principal ideal $(t^N)$ in $\mathbb Q[t]$.
Consequently a section exists exactly when $c\in\mathbb Q$.
The pullback assertion follows from the defining equation.
The idempotent $e_p$ has value one at $p$ and zero at every
other prime; it cannot be a diagonal rational number, since
a rational number with a zero image in one $\mathbb Q_\ell$
is zero. This proves the final claim.
\end{proof}

\paragraph{AT04 (Proved; the positive source-weight comparison).}
For $u\geq0$ and $N\geq1$, restriction along
$I_N\langle u\rangle\hookrightarrow S\langle u\rangle$
induces
\[
 R\operatorname{Hom}_{\mathrm{gr}S}
       (S\langle u\rangle,S)
 \xrightarrow{\ \cong\ }
 R\operatorname{Hom}_{\mathrm{gr}S}
       (I_N\langle u\rangle,S).
\]
Both sides are $S^u$ in cohomological degree zero and zero
in every other degree.

\begin{proof}
Shift both source and target by $-u$. The right side has
the Ext groups of $I_N$ into $S\langle-u\rangle$.
Since $-u\leq0$, AT02 makes its positive Ext groups zero
and identifies its degree-zero Hom by restriction from
$S$. The left side has exactly these same groups because
its source is free. The indicated comparison therefore
induces an isomorphism on every cohomology group.
\end{proof}

AT04 applies to the weight direction of a delayed
positive-weight cell. AT03 proves persistence in the
opposite shift direction. We have not identified either
calculation with an invariant of admissible theories.
In particular we do not infer an actual simplicial
PD-algebra primitive from a derived module Hom computation:
the coefficient-tail parameter is not a free $S$-module,
and that strictification step still requires a construction.

\paragraph{AT05 (Proved; degree-preserving orthogonality of nonnegative torsion).}
Let $M$ be any nonnegatively weight-graded $S$-module.
Let $T$ be any such module whose every element is killed
by all sufficiently high coefficient basis elements.
Then
\[
 R\operatorname{Hom}_{\mathrm{gr}S}(T,S)=0.
\]
In particular for every $N\geq0$, restriction gives an
isomorphism
\[
 R\operatorname{Hom}_{\mathrm{gr}S}(M,S)
 \xrightarrow{\ \cong\ }
 R\operatorname{Hom}_{\mathrm{gr}S}(M_{\geq N},S).
\]
These are degree-preserving Hom complexes. No assertion
of the vanishing of all internal weight shifts is made.

\begin{proof}
Over each $S_p$, the flat base change $T_p$ is a torsion
module: every element is a finite sum of simple tensors,
so a common annihilating coefficient tail can be chosen.
Theorem~8.3 of Nagpal--Snowden applies with finitely
presented target $S_p$, and gives zero derived Hom.
Derived adjunction and the graded product argument of
AT02 give
\[
 R\operatorname{Hom}_{\mathrm{gr}S}
       (T,S_{\widehat{\mathbb Z}})=0.
\]

We verify the rational terms separately. Every finitely
generated homogeneous torsion submodule of $T$ is bounded
in weight: choose a coefficient-tail annihilator for each
of its finitely many homogeneous generators and take the
maximum of the resulting upper weight bounds. Its
rationalization is a finite-dimensional nonnegatively
graded $\mathbb Q[t]$-module. Filter it by the finitely many
weight tails. Each successive quotient is a finite sum
of copies of $\mathbb Q\langle r\rangle$ with $r\geq0$.
For any rational coefficient vector space $V$, the graded
free resolution
\[
 0\longrightarrow\mathbb Q[t]\langle r+1\rangle
 \xrightarrow{t}\mathbb Q[t]\langle r\rangle
 \longrightarrow\mathbb Q\langle r\rangle\longrightarrow0
\]
has degree-preserving Hom into $V[t]$ equal to
$Vt^r\xrightarrow{t}Vt^{r+1}$, an isomorphism.
Thus all these derived Hom complexes vanish; the finite
filtration gives the same conclusion for the finitely
generated torsion submodule.

The module $T$ is the filtered union of those submodules.
Filtered colimits of modules are exact, so this union is
also their homotopy colimit in the derived category.
Derived Hom out of it is the homotopy limit of their
derived Hom complexes. Each complex is zero, so the
homotopy limit is zero. Taking $V=\mathbb Q$ and
$V=\mathbb A_f$, and using flat base change as in AT02,
shows that both rational terms of the arithmetic sequence
have zero derived Hom from $T$. Applying derived Hom to
AT01 now proves the first assertion over $\mathbb Z$.

Finally $M/M_{\geq N}$ is bounded in nonnegative weights,
and is therefore torsion. Apply derived Hom into $S$ to
$0\to M_{\geq N}\to M\to M/M_{\geq N}\to0$.
The just-proved vanishing gives the asserted restriction
isomorphism.
\end{proof}
